Recall the results we have proven
\begin{enumerate}
	\item In the simplest Semi-Online model, Ranking has a competitiveness of $1-\frac{1}{e}=0.632$ and is optimal \label{RNK}
	\item If vertices arrive in a random order, Ranking has a competitiveness above $0.696$ \label{RA1}
	\item If vertices arrive in a random order, any deterministic algorithm with consistent tie-breaking has a competitiveness of $1-\frac{1}{e}$ \label{RA2}
	\item If vertices arrive from a known distribution, a deterministic TSM algorithm has a competitiveness of $0.670$ \label{TSM}
	\item Even a substantial delay in decision time doesn't increase competitiveness \label{delay}
	\item Prioritizing matching based on the knowledge of final degrees doesn't increase competitiveness and may even reduce it for certain randomized algorithms \label{RED}
	\item In a more random graph the usage of final degrees leads to an optimal algorithm (MPD) \label{MPD}
	\item A high recourse cannot be forced \label{Recourse}
	\item Waterlevel fractional matching has a competitiveness of $1-\frac{1}{e}=0.632$ \label{FRAC}
	\item A deterministic algorithm, even allowed to make small corrections, has a good competitiveness of $\frac{d+1}{d+2}$, being more competitive than Ranking in the Semi-Online model with just $1$ correction per arriving vertex. \label{COR}
	\item none of the mentioned competitiveness models allow for a perfect algorithm, which always finds a perfect matching, ie with a competitiveness ratio of $1$.
\end{enumerate}

and our research questions

\begin{enumerate}
	\item what modifications are required for deterministic algorithms to break the $0.5$ competitiveness barrier
	\item is there a model where a perfect algorithm which always finds a maximum matching is possible
\end{enumerate}

In terms of breaking the $0.5$ deterministic competitiveness barrier, we present a variety of ways of accomplishing this, changing the arrival to random or stochastic instead of adversarial, allowing for corrections (even as small as $1$ edge per arrival) and allowing for non-commitment by finding a fractional matching instead all allow deterministic algorithms to achieve competitiveness above $0.5$. What is also interesting is that neither knowledge of final degrees nor knowing up to half of all online vertices in advance allow for this. The pattern seems to be that only providing information on the final graph, without constraining the adversary, is not sufficient for this task. It would seem that in this setting constraining the adversary is more important than providing information for better choices. Indeed, using this additional information might be detrimental, as in the case of the Final Degree model.

Another interesting pattern is in whether randomness is utilized in various models. In random arrivals, randomized Ranking performs better than deterministic algorithms with consistent tie-breaking while in stochastic matching, where more information is available, considered algorithms are deterministic. It would appear that, when the adversary is constrained, the ability to make better choices makes randomness less valuable.

We also note that, in the case of Online Bipartite Matching with One-Sided Vertex Arrivals, allowing for corrections appears to be a very strong algorithm enhancement as the required correction budges to surpass Ranking in the Semi-Online Model in terms of competitiveness is only $d=1$ yielding an algorithm with a competitiveness of $\frac{2}{3}>1-\frac{1}{e}$ and maintaining a significantly lower Recourse of $O(n\log^2n)$ is achievable compared to Edge Arrival or Two-Sided Vertex Arrival.

Lastly, none of these models, aside from Recourse where competitiveness isn't a concern, allows for a perfect algorithm. It is unsurprising as we know that a matching is maximum iff there are no augmenting paths and no models we showcased, and indeed no conventional modification to models of online problems, gives this information. It is therefore unsurprising that making the right decision as a vertex arrives is difficult.